\IfFileExists{stacks-project.cls}{%
\documentclass{stacks-project}
}{%
\documentclass{amsart}
}

% For dealing with references we use the comment environment
\usepackage{verbatim}
\newenvironment{reference}{\comment}{\endcomment}
%\newenvironment{reference}{}{}
\newenvironment{slogan}{\comment}{\endcomment}
\newenvironment{history}{\comment}{\endcomment}

% For commutative diagrams we use Xy-pic
\usepackage[all]{xy}

% We use 2cell for 2-commutative diagrams.
\xyoption{2cell}
\UseAllTwocells

% We use multicol for the list of chapters between chapters
\usepackage{multicol}

% This is generally recommended for better output
\usepackage{lmodern}
\usepackage[T1]{fontenc}

% For cross-file-references
\usepackage{xr-hyper}

% Package for hypertext links:
\usepackage{hyperref}

% For any local file, say "hello.tex" you want to link to please
% use \externaldocument[hello-]{hello}
\externaldocument[introduction-]{introduction}
\externaldocument[conventions-]{conventions}
\externaldocument[sets-]{sets}
\externaldocument[categories-]{categories}
\externaldocument[topology-]{topology}
\externaldocument[sheaves-]{sheaves}
\externaldocument[sites-]{sites}
\externaldocument[stacks-]{stacks}
\externaldocument[fields-]{fields}
\externaldocument[algebra-]{algebra}
\externaldocument[brauer-]{brauer}
\externaldocument[homology-]{homology}
\externaldocument[derived-]{derived}
\externaldocument[simplicial-]{simplicial}
\externaldocument[more-algebra-]{more-algebra}
\externaldocument[smoothing-]{smoothing}
\externaldocument[modules-]{modules}
\externaldocument[sites-modules-]{sites-modules}
\externaldocument[injectives-]{injectives}
\externaldocument[cohomology-]{cohomology}
\externaldocument[sites-cohomology-]{sites-cohomology}
\externaldocument[dga-]{dga}
\externaldocument[dpa-]{dpa}
\externaldocument[sdga-]{sdga}
\externaldocument[hypercovering-]{hypercovering}
\externaldocument[schemes-]{schemes}
\externaldocument[constructions-]{constructions}
\externaldocument[properties-]{properties}
\externaldocument[morphisms-]{morphisms}
\externaldocument[coherent-]{coherent}
\externaldocument[divisors-]{divisors}
\externaldocument[limits-]{limits}
\externaldocument[varieties-]{varieties}
\externaldocument[topologies-]{topologies}
\externaldocument[descent-]{descent}
\externaldocument[perfect-]{perfect}
\externaldocument[more-morphisms-]{more-morphisms}
\externaldocument[flat-]{flat}
\externaldocument[groupoids-]{groupoids}
\externaldocument[more-groupoids-]{more-groupoids}
\externaldocument[etale-]{etale}
\externaldocument[chow-]{chow}
\externaldocument[intersection-]{intersection}
\externaldocument[pic-]{pic}
\externaldocument[weil-]{weil}
\externaldocument[adequate-]{adequate}
\externaldocument[dualizing-]{dualizing}
\externaldocument[duality-]{duality}
\externaldocument[discriminant-]{discriminant}
\externaldocument[derham-]{derham}
\externaldocument[local-cohomology-]{local-cohomology}
\externaldocument[algebraization-]{algebraization}
\externaldocument[curves-]{curves}
\externaldocument[resolve-]{resolve}
\externaldocument[models-]{models}
\externaldocument[functors-]{functors}
\externaldocument[equiv-]{equiv}
\externaldocument[pione-]{pione}
\externaldocument[etale-cohomology-]{etale-cohomology}
\externaldocument[proetale-]{proetale}
\externaldocument[relative-cycles-]{relative-cycles}
\externaldocument[more-etale-]{more-etale}
\externaldocument[trace-]{trace}
\externaldocument[crystalline-]{crystalline}
\externaldocument[spaces-]{spaces}
\externaldocument[spaces-properties-]{spaces-properties}
\externaldocument[spaces-morphisms-]{spaces-morphisms}
\externaldocument[decent-spaces-]{decent-spaces}
\externaldocument[spaces-cohomology-]{spaces-cohomology}
\externaldocument[spaces-limits-]{spaces-limits}
\externaldocument[spaces-divisors-]{spaces-divisors}
\externaldocument[spaces-over-fields-]{spaces-over-fields}
\externaldocument[spaces-topologies-]{spaces-topologies}
\externaldocument[spaces-descent-]{spaces-descent}
\externaldocument[spaces-perfect-]{spaces-perfect}
\externaldocument[spaces-more-morphisms-]{spaces-more-morphisms}
\externaldocument[spaces-flat-]{spaces-flat}
\externaldocument[spaces-groupoids-]{spaces-groupoids}
\externaldocument[spaces-more-groupoids-]{spaces-more-groupoids}
\externaldocument[bootstrap-]{bootstrap}
\externaldocument[spaces-pushouts-]{spaces-pushouts}
\externaldocument[spaces-chow-]{spaces-chow}
\externaldocument[groupoids-quotients-]{groupoids-quotients}
\externaldocument[spaces-more-cohomology-]{spaces-more-cohomology}
\externaldocument[spaces-simplicial-]{spaces-simplicial}
\externaldocument[spaces-duality-]{spaces-duality}
\externaldocument[formal-spaces-]{formal-spaces}
\externaldocument[restricted-]{restricted}
\externaldocument[spaces-resolve-]{spaces-resolve}
\externaldocument[formal-defos-]{formal-defos}
\externaldocument[defos-]{defos}
\externaldocument[cotangent-]{cotangent}
\externaldocument[examples-defos-]{examples-defos}
\externaldocument[algebraic-]{algebraic}
\externaldocument[examples-stacks-]{examples-stacks}
\externaldocument[stacks-sheaves-]{stacks-sheaves}
\externaldocument[criteria-]{criteria}
\externaldocument[artin-]{artin}
\externaldocument[quot-]{quot}
\externaldocument[stacks-properties-]{stacks-properties}
\externaldocument[stacks-morphisms-]{stacks-morphisms}
\externaldocument[stacks-limits-]{stacks-limits}
\externaldocument[stacks-cohomology-]{stacks-cohomology}
\externaldocument[stacks-perfect-]{stacks-perfect}
\externaldocument[stacks-introduction-]{stacks-introduction}
\externaldocument[stacks-more-morphisms-]{stacks-more-morphisms}
\externaldocument[stacks-geometry-]{stacks-geometry}
\externaldocument[moduli-]{moduli}
\externaldocument[moduli-curves-]{moduli-curves}
\externaldocument[examples-]{examples}
\externaldocument[exercises-]{exercises}
\externaldocument[guide-]{guide}
\externaldocument[desirables-]{desirables}
\externaldocument[coding-]{coding}
\externaldocument[obsolete-]{obsolete}
\externaldocument[fdl-]{fdl}
\externaldocument[index-]{index}

% Theorem environments.
%
\theoremstyle{plain}
\newtheorem{theorem}[subsection]{Theorem}
\newtheorem{proposition}[subsection]{Proposition}
\newtheorem{lemma}[subsection]{Lemma}

\theoremstyle{definition}
\newtheorem{definition}[subsection]{Definition}
\newtheorem{example}[subsection]{Example}
\newtheorem{exercise}[subsection]{Exercise}
\newtheorem{situation}[subsection]{Situation}

\theoremstyle{remark}
\newtheorem{remark}[subsection]{Remark}
\newtheorem{remarks}[subsection]{Remarks}

\numberwithin{equation}{subsection}

% Macros
%
\def\lim{\mathop{\mathrm{lim}}\nolimits}
\def\colim{\mathop{\mathrm{colim}}\nolimits}
\def\Spec{\mathop{\mathrm{Spec}}}
\def\Hom{\mathop{\mathrm{Hom}}\nolimits}
\def\Ext{\mathop{\mathrm{Ext}}\nolimits}
\def\SheafHom{\mathop{\mathcal{H}\!\mathit{om}}\nolimits}
\def\SheafExt{\mathop{\mathcal{E}\!\mathit{xt}}\nolimits}
\def\Sch{\mathit{Sch}}
\def\Mor{\mathop{\mathrm{Mor}}\nolimits}
\def\Ob{\mathop{\mathrm{Ob}}\nolimits}
\def\Sh{\mathop{\mathit{Sh}}\nolimits}
\def\NL{\mathop{N\!L}\nolimits}
\def\CH{\mathop{\mathrm{CH}}\nolimits}
\def\proetale{{pro\text{-}\acute{e}tale}}
\def\etale{{\acute{e}tale}}
\def\QCoh{\mathit{QCoh}}
\def\Ker{\mathop{\mathrm{Ker}}}
\def\Im{\mathop{\mathrm{Im}}}
\def\Coker{\mathop{\mathrm{Coker}}}
\def\Coim{\mathop{\mathrm{Coim}}}

% Boxtimes
%
\DeclareMathSymbol{\boxtimes}{\mathbin}{AMSa}{"02}

%
% Macros for moduli stacks/spaces
%
\def\QCohstack{\mathcal{QC}\!\mathit{oh}}
\def\Cohstack{\mathcal{C}\!\mathit{oh}}
\def\Spacesstack{\mathcal{S}\!\mathit{paces}}
\def\Quotfunctor{\mathrm{Quot}}
\def\Hilbfunctor{\mathrm{Hilb}}
\def\Curvesstack{\mathcal{C}\!\mathit{urves}}
\def\Polarizedstack{\mathcal{P}\!\mathit{olarized}}
\def\Complexesstack{\mathcal{C}\!\mathit{omplexes}}
% \Pic is the operator that assigns to X its picard group, usage \Pic(X)
% \Picardstack_{X/B} denotes the Picard stack of X over B
% \Picardfunctor_{X/B} denotes the Picard functor of X over B
\def\Pic{\mathop{\mathrm{Pic}}\nolimits}
\def\Picardstack{\mathcal{P}\!\mathit{ic}}
\def\Picardfunctor{\mathrm{Pic}}
\def\Deformationcategory{\mathcal{D}\!\mathit{ef}}
\usepackage{cleveref}
\crefname{originallemma}{original lemma}{original lemmas}
\Crefname{originallemma}{Original lemma}{Original lemmas}
\makeatletter
\newlabel{lemma-equivalence-quasi-coherent-limits}{{06VE}{source}{}{sourcepointer-1}{}}
\newlabel{lemma-equivalence-quasi-coherent-limits@cref}{{[originallemma][0][]06VE}{[1][0][]source}{}{}{}}
\makeatother
\crefname{originallemma}{original lemma}{original lemmas}
\Crefname{originallemma}{Original lemma}{Original lemmas}
\makeatletter
\newlabel{lemma-equivalence-quasi-coherent-properties}{{05VG}{source}{}{sourcepointer-2}{}}
\newlabel{lemma-equivalence-quasi-coherent-properties@cref}{{[originallemma][0][]05VG}{[1][0][]source}{}{}{}}
\makeatother
\crefname{originalsection}{original section}{original sections}
\Crefname{originalsection}{Original section}{Original sections}
\makeatletter
\newlabel{schemes-section-quasi-coherent}{{01LA}{source}{}{sourcepointer-3}{}}
\newlabel{schemes-section-quasi-coherent@cref}{{[originalsection][0][]01LA}{[1][0][]source}{}{}{}}
\makeatother
\title{Properties of the category of quasi-coherent modules on the small etale site}
\author{The Stacks Project Authors}
\date{stacks-project:descent.tex:local-snapshot-20261012; GNU FDL 1.2 or later}
\begin{document}
\maketitle
\noindent\textit{Editorial note.} All delivered mathematical text is the author's own, copied from the pinned Stacks Project chapter file descent.tex; no statement or proof step was rewritten, reordered or removed.\par\medskip
\noindent\textit{Editorial note.} The counted spans carry navigation stubs for the references leaving them; no bibliography citation occurs in the spans.\par\medskip
\noindent\textit{Editorial note.} The preamble span points at the chapter's own preamble.tex shipped with the frozen source, used verbatim as staged.\par\medskip
\noindent\textit{Allowed standard prerequisites (editorial).}\begin{itemize}
\item Compatibility of the equivalence with limits, colimits, kernels, cokernels, tensor products and internal homs. Same chapter, Lemma lemma-equivalence-quasi-coherent-limits, invoked at proof lines 3453-3454 and at each part. Carries each of the six properties across the equivalence between quasi-coherent modules on S and on S\_etale.
\item Preservation of finite presentability along the equivalence. Same chapter, Lemma lemma-equivalence-quasi-coherent-properties, invoked at proof lines 3505-3506. Used in part (6) to transfer the finite presentability hypothesis.
\item The corresponding closure properties for quasi-coherent modules on a scheme. Schemes, Section schemes-section-quasi-coherent, invoked at proof lines 3451-3452. Supplies the facts on the scheme side that are being transferred.
\end{itemize}\medskip
\ifdefined\hypertarget\hypertarget{sourcepointer-1}{}\fi
\noindent\textit{Source reference.} Original lemma 06VE: sources/stacks-descent/pinned-original-source.tex, source lines 3256--3256. Pointer outside the counted span, invoked once for context; the load-bearing argument of the counted statement is the author's proof reproduced in full in the proof span.\par\medskip
\ifdefined\hypertarget\hypertarget{sourcepointer-2}{}\fi
\noindent\textit{Source reference.} Original lemma 05VG: sources/stacks-descent/pinned-original-source.tex, source lines 2830--2830. Pointer outside the counted span, invoked once for context; the load-bearing argument of the counted statement is the author's proof reproduced in full in the proof span.\par\medskip
\ifdefined\hypertarget\hypertarget{sourcepointer-3}{}\fi
\noindent\textit{Source reference.} Original section 01LA: sources/stacks-chow/tags.txt, source lines 2079--2079. Pointer outside the counted span, invoked once for context; the load-bearing argument of the counted statement is the author's proof reproduced in full in the proof span.\par\medskip
\begin{lemma}
\label{lemma-properties-quasi-coherent}
Let $S$ be a scheme. The category $\QCoh(S_\etale, \mathcal{O})$
of quasi-coherent modules on $S_\etale$
has the following properties:
\begin{enumerate}
\item Any direct sum of quasi-coherent sheaves is quasi-coherent.
\item Any colimit of quasi-coherent sheaves is quasi-coherent.
\item The kernel and cokernel of a morphism of quasi-coherent sheaves
is quasi-coherent.
\item Given a short exact sequence of $\mathcal{O}$-modules
$0 \to \mathcal{F}_1 \to \mathcal{F}_2 \to \mathcal{F}_3 \to 0$
if two out of three are quasi-coherent so is the third.
\item Given two quasi-coherent $\mathcal{O}$-modules
the tensor product is quasi-coherent.
\item Given two quasi-coherent $\mathcal{O}$-modules
$\mathcal{F}$, $\mathcal{G}$ such that $\mathcal{F}$
is of finite presentation.
then the internal hom
$\SheafHom_\mathcal{O}(\mathcal{F}, \mathcal{G})$
is quasi-coherent.
\end{enumerate}
\end{lemma}

\begin{proof}
The corresponding facts hold for quasi-coherent modules on the scheme $S$,
see Schemes, Section \ref{schemes-section-quasi-coherent}. The proof will
be to use Lemma \ref{lemma-equivalence-quasi-coherent-limits} to transfer
these truths to $S_\etale$.

\medskip\noindent
Proof of (1). Let $\mathcal{F}_i$, $i \in I$ be a family of objects of
$\QCoh(S_\etale, \mathcal{O})$. Write $\mathcal{F}_i = \mathcal{G}_i^a$
for some quasi-coherent modules $\mathcal{G}_i$ on $S$.
Then $\bigoplus \mathcal{F}_i = (\bigoplus \mathcal{G}_i)^a$ by
the lemma cited and we conclude.

\medskip\noindent
Proof of (2). Let $\mathcal{I} \to \QCoh(S_\etale, \mathcal{O})$,
$i \mapsto \mathcal{F}_i$ be a diagram. Write
$\mathcal{F}_i = \mathcal{G}_i^a$ so we get a diagram
$\mathcal{I} \to \QCoh(\mathcal{O}_S)$.
Then $\colim \mathcal{F}_i = (\colim \mathcal{G}_i)^a$ by
the lemma cited and we conclude.

\medskip\noindent
Proof of (3). Let $a : \mathcal{F} \to \mathcal{F}'$
be an arrow of $\QCoh(S_\etale, \mathcal{O})$.
Write $a = b^a$ for some map $b : \mathcal{G} \to \mathcal{G}'$
of quasi-coherent modules on $S$. By the lemma cited
we have $\Ker(a) = \Ker(b)^a$ and $\Coker(a) = \Coker(b)^a$
and we conclude.

\medskip\noindent
Proof of (4). This follows from (3) except in the case when
we know $\mathcal{F}_1$ and $\mathcal{F}_3$ are quasi-coherent.
In this case write $\mathcal{F}_1 = \mathcal{G}_1^a$
and $\mathcal{F}_3 = \mathcal{G}_3^a$ with
$\mathcal{G}_i$ quasi-coherent on $S$.
By Lemma \ref{lemma-equivalence-quasi-coherent-limits} part (10)
we conclude.

\medskip\noindent
Proof of (5). Let $\mathcal{F}$ and $\mathcal{F}'$
be in $\QCoh(S_\etale, \mathcal{O})$.
Write $\mathcal{F} = \mathcal{G}^a$
and $\mathcal{F}' = (\mathcal{G}')^a$
with $\mathcal{G}$ and $\mathcal{G}'$ quasi-coherent on $S$.
By the lemma cited we have
$\mathcal{F} \otimes_\mathcal{O} \mathcal{F}' =
(\mathcal{G} \otimes_{\mathcal{O}_S} \mathcal{G}')^a$
and we conclude.

\medskip\noindent
Proof of (6). Let $\mathcal{F}$ and $\mathcal{G}$
be in $\QCoh(S_\etale, \mathcal{O})$ with $\mathcal{F}$
of finite presentation. Write $\mathcal{F} = \mathcal{H}^a$
and $\mathcal{G} = (\mathcal{I})^a$
with $\mathcal{H}$ and $\mathcal{I}$ quasi-coherent on $S$.
By Lemma \ref{lemma-equivalence-quasi-coherent-properties}
we see that $\mathcal{H}$ is of finite presentation.
By Lemma \ref{lemma-equivalence-quasi-coherent-limits} part (8)
we have
$\SheafHom_\mathcal{O}(\mathcal{F}, \mathcal{G}) =
(\SheafHom_{\mathcal{O}_S}(\mathcal{H}, \mathcal{I}))^a$
and we conclude.
\end{proof}

\end{document}
